\subsubsection*{The theta polynomials from pair removal}
\hypertarget{source-twist-euler}{}

The toric distribution relation describes the effect of adding a split prime. To use it in the simultaneous choice of Section~\ref{subsec:retained-twist}, encode pair removal by specialization in the group ring before evaluating any character.
The first and second halves use the multiplicative group-ring coordinates
\[
 U_i=1+X_i,\qquad V_i=1+Y_i,\qquad U_i^L=V_i^L=1,\qquad L=\ell^M.
\]

The group ring of the coordinates is the quotient by
\(U_i^L-1,V_i^L-1\). In an integral coefficient polynomial one
can use the exponent \(L-1\) for the inverse of a group
coordinate. Equality at all \(L\)-torsion tuples after setting
\(V_i=1\) gives equality modulo \((Y_i,U_i^L-1)_i\) over the
original integers. Indeed division in each variable by the
monic polynomial leaves a remainder of degree less than \(L\)
in each variable. In a characteristic-zero splitting field
successive univariate interpolation kills that remainder.
This proof involves no division by a group order in the
coefficient ring.

\begin{lemma}[Removal of pairs for abelian theta values]\label{lem:abelian-pair-removal}
For an abelian label at an admissible critical integer \(s\ge0\),
use the uniform denominator \(\ell^{N_0}\) and the two halves of the
chosen auxiliary pairs. The cleared full and first-half theta values
have integral polynomial representatives \(f_M,g_M\) such that
\[
 f_M\equiv B_Mg_M\pmod{(Y_i,U_i^L-1)_i},\qquad
 v(g_M(0))=N_0+v(\theta_\lambda(1)).
\]
Each second-half factor of \(B_M\) reduces to a polynomial in
\(X_i\) of order zero or two at zero.
\end{lemma}

\begin{proof}
First consider an abelian label at an admissible critical
integer \(s\ge0\). The input denominator \(\ell^{N_0}\)
makes each coefficient of its ray-class theta element integral,
uniformly in the auxiliary level. Push this element to the
coordinate group ring by the chosen characters and choose a
polynomial lift \(f_M\). More explicitly, lift each ray class
to an Artin element, use its character coordinates as exponents
in \([0,L-1]\), and sum the resulting monomials in \(U_i,V_i\)
with the cleared coefficients. The characters factor through
the ray group. Evaluating gives exactly
\(\ell^{N_0}\theta_\lambda(\psi)\). Use \(g_M\) for the
same construction at the union of the first half of the pairs,
using just the first-half characters.

\proofstep{The second-half factors.}
The pair-removal identity of the theta input at a pair of
norm \(q\) and at a character unramified on that pair gives
the factor
\[
 E_q(A)=(1-Aq^{-s})(1-A^{-1}q^{-s})+(q-1)q^{-2s},
 \qquad A=\lambda(\mathrm{Frob})\psi(\mathrm{Frob}).
\]
All these coefficients are integral at \(\ell\) because
\(q\equiv1\bmod 2L\). Remove the second half of the pairs
with all \(V_i=1\). At pair \(i\) the Frobenius character
value is exactly \(U_i\) on evaluating, by the identity
matrix condition on the first-half characters. Thus \(B_M\)
is the product over the second half of the integral polynomial
representatives of \(E_q(\lambda(\mathrm{Frob})U_i)\).
Successive removal and the interpolation argument above give
\[
 f_M\equiv B_Mg_M\bmod(Y_i,U_i^L-1)_i.
\]


\proofstep{The first-half unit test.}
To compute \(g_M(0)\), remove the first half at the trivial
character. Choose the open normal subgroup inside the kernel
of the residue character \(\bar\lambda\). On the prescribed
tame element \(z\), \(\bar\lambda(z)\ne1\). Indeed at the
distinguished place of norm \(q_0\) the character label there
agrees with \(\delta^{\pm1}\) after taking a \(q_0\)-power
iterate (power \(q_0^n\) for some \(n\)). The finite
quotient values of \(z\) have order prime to \(q_0\ell\).
Writing a bar for reduction at \(\ell\), the choice of \(z\) gives
\[
 (\overline{\delta(z)}-1)(\overline{\delta(z)}+1)=\overline{\delta(z)}^2-1\ne0.
\]
These are the tests on \(z\): take the
prime-to-\(q_0\ell\) part of a generator at the distinguished
order-15 place. Its value retains the prime-to-\(\ell\)
part of a primitive fifteenth root, of order divisible by
three or five and prime to \(q_0\). It follows that
\(\bar\lambda(z)^{q_0^n}\ne1\).

By the first Frobenius test, the factors at the first-half
places reduce at the trivial character to
\((1-\bar\lambda(z))(1-\bar\lambda(z)^{-1})\), a unit.
Consequently
\[
 v(g_M(0))=N_0+v(\theta_\lambda(1)).
\]
Here the value on the right is at empty auxiliary level.
For a second-half factor write \(x=\bar\lambda(\mathrm{Frob})\).
Its reduction is \((1-xU_i)(1-x^{-1}U_i^{L-1})\).
For \(x\ne1\) it has unit constant term. For \(x=1\)
it equals \(X_i^2\sum_{j=0}^{L-2}(1+X_i)^j\), whose last
factor has constant \(L-1\ne0\) in the residue field.
The denominator bound and the pair relation prove the abelian coefficient condition with \(n_*=2\). The distinction between the
two halves of the primes is essential to the constant
coefficient estimate. \label{cor:twist-euler-1}
\end{proof}

\begin{proposition}[Determinant and signed index for the row theta values]\label{prop:row-pair-removal}
Let \(B\) be the prescribed row lattice of rank \(n\), with its
\(n\) independent projected vectors at empty level and the uniform
clearing scalar \(x\). Write \(C_Q\) for their cleared coordinate
matrix at level \(Q\), \(Q_1\) for the first half of the auxiliary
pairs, and \(b=\operatorname{ord}_O(\det C_0)\). Then
\(f_M=\det C_Q\) and \(g_M=\det C_{Q_1}\) satisfy the same
pair-removal congruence as in Lemma~\ref{lem:abelian-pair-removal},
with the second-half determinant factors displayed below, of residual
order at most \(rn\) in each variable. For any lattice containing the specialized columns, let \(\Delta\) denote its excess signed index over the original column index. Then
\[
 \operatorname{ord}_O g_M(0)=b,\qquad f_M(u-1)\ne0\ \Longrightarrow\ \Delta\le v(f_M(u-1))-b/e_O.
\]
\end{proposition}

\begin{proof}
\proofstep{Clearing the columns on both sides.}
For a row, let \(B\) be the prescribed free lattice of rank
\(n\), \(A\) the commuting central operator of the row
calculation, and \(\theta_1,\ldots,\theta_n\) the selected
projected vectors at empty level. They are independent and
span the rational space. Choose a basis of \(B\).
There is a single denominator
\[
 x=\ell^c\ne0
\]
for the
universal lifts of this finite list, valid at every \(M,Q\)
disjoint from their levels. Projection takes the lattice
containing these cleared lifts into \(B\).

Take their basis
coordinates, with the universal characters still present,
and push the group coordinates to the polynomial presentation
above. This gives an integral square matrix \(C_Q\) such
that at every torsion character its columns are the coordinates
of \(x\) times the projected specializations. At empty level
it gives the constant matrix \(C_0\) of the \(x\theta_j\).
In particular \(\det C_0\ne0\). With \(Q_1\) the first-half set, the two determinant polynomials and the base valuation are
\[
 f_M=\det C_Q,\qquad g_M=\det C_{Q_1},\qquad b=\operatorname{ord}_O(\det C_0).
\]

This choice cancels the denominator in the signed index.
Indeed, if \(\det C_Q(u-1)\ne0\), those specialized columns
before clearing span a lattice of normalized relative index
\(v(\det C_Q(u-1))-n v(x)\) in \(B\). Every lattice in
the same rational space that contains them has relative
index at most this, by elementary divisors. Subtract the
corresponding index for the original \(\theta_j\) to get
exactly \(v(f_M(u-1))-b/e_O\). The formula also holds for
fractional lattices extending beyond \(B\).

\proofstep{The primitive local function.}
For a place \(P\) of norm \(q\), with integral Hecke matrices
\(T_{a,P}\) in the basis of \(B\), put
\[
 H_P(Y)=I+\sum_{a=1}^r(-1)^a q^{a(a-1)/2}Y^aT_{a,P}.
\]
Work in the split local model at the chosen member \(P\).
To distinguish the matrix variable from the rank of the row
lattice, denote its number of columns by \(n_{\rm loc}\).
It has \(r\) rows. The split base frames
identify the vacuum with the characteristic function of
integral matrices. Sum over the character coordinate at a
deleted pair. If the reduced rows span an \(r\)-plane, its
complements are the graphs of maps from a fixed complement
to that plane. For a matrix \(z\) the normalized count is
\[
 q^{-r(n_{\rm loc}-r)}
 \#\{\text{complements of dimension } n_{\rm loc}-r\}
 =
 \begin{cases}1,&\operatorname{rk}\bar z=r,\\0,&\operatorname{rk}\bar z<r.\end{cases}
\]
For smaller rank the dimension condition prevents complementarity.
This is the primitive local function on the integral vacuum.

Let \(z\) be integral and \(c=\dim\ker(\bar z^{\mathsf T})\).
The right cosets of the index-\(a\) minuscule Hecke double
coset parameterize the lattices \(D\) between
\(\varpi O_P^r\) and \(O_P^r\) of residual codimension \(a\):
send a representative \(\gamma\) to \(\gamma O_P^r\), and
use elementary divisors for surjectivity. The condition
\(\gamma^{-1}z\) integral means that each column is in
\(D\). Equivalently the annihilator of \(D/\varpi O_P^r\)
is an \(a\)-plane in \(\ker(\bar z^{\mathsf T})\).
Thus there are \(\genfrac{[}{]}{0pt}{}{c}{a}_q\) translates.
The subspace recursion obtained by intersection with a hyperplane
gives the polynomial identity
\[
 \sum_{a=0}^c (-1)^a q^{a(a-1)/2}
    \genfrac{[}{]}{0pt}{}{c}{a}_q Z^a
   =\prod_{j=0}^{c-1}(1-q^jZ).
\]
At \(Z=1\) the product is one for \(c=0\) and zero for \(c>0\).
Hence the alternating Hecke sum is precisely the primitive
indicator. Off the integral vacuum every term vanishes since
\(D\subset O_P^r\).

\proofstep{The normalized Hecke translates.}
In the split Weil action the right translate of the vacuum
by an index-\(a\) representative has the indicator just
counted with modulus and splitting-character multiplier
\[
 q^{a n_{\rm loc}/2}\mu(\varpi)^a.
\]
Here \(n_{\rm loc}\) is even. This is
the change-of-variables factor \(|\det\gamma|^{-n_{\rm loc}/2}\)
on the matrix argument and the unramified splitting
character on its determinant. Dividing the sum of translates
by this same multiplier is the normalized \(T_{a,P}\) on the
row.

The other local components carry the shorter coordinate
datum: the coordinates agree by restriction, and replacing
the deleted test by the constant one on its integral support
restores the vacuum. The identity of local functions passes
to theta series by linearity and right translation. If \(b\) is the remaining Artin label, the determinant of a representative in degree \(a\) changes it by
\[
 b\longmapsto b+a\delta,
\]
where \(\delta\) is the Artin-coordinate vector of the uniformizer at \(P\). This follows
by taking determinants of the representatives, all of which
are \(\varpi^a\) times a local unit and all remaining
characters are unramified at \(P\). With the inverse Artin map of Lemma~\ref{lem:toric-split-distribution} \(\delta\) is minus the vector of
their arithmetic Frobenius values.

For a character descending to the shorter set, write \(Y\) for its Frobenius value. Translating the index back by \(-a\delta\) supplies the factor
\[
 \chi(-a\delta)=\chi(-\delta)^a=Y^a.
\]
Injectivity of
the row realization transfers this equality to the prescribed
Hecke matrices. This gives \(H_P(Y)\) at the signed character
value on the representative selected by the row.

Projection commutes with the
Hecke operators. Each of the \(n\) columns therefore changes
by the very same matrix. Applying the relation successively
and taking determinants gives
\[
 B_M=\prod_{i=1}^{R_*} \det H_{P'_i}(U_i^{c_i}),
       \qquad c_i\in\{1,L-1\},
\]
with \(P'_i\) in the second half. The identity-matrix
Frobenius test gives \(U_i\) on the specified representative,
conjugating the place inverts the value by anticyclotomicity,
and the sign of the row can invert it once more. This proves
the formula at torsion tuples with trivial second half and
hence the required ideal congruence by interpolation.

Over the residue field
\[
 h(Y)=\det H_{P'_i}(Y)
\]
satisfies
\[
 h(0)=1,\qquad\deg h\le rn.
\] The order of
\(h((1+X)^{c_i})\) at \(X=0\) is the multiplicity of the
root \(Y=1\) of \(h\), including multiplicity zero. Indeed
\((1+X)^{c_i}-1\) has nonzero linear coefficient there.
Each factor is therefore \(X_i^{n_i}\) times a polynomial
with unit constant, \(n_i\le rn\), uniformly in \(M\).

\proofstep{Units at the first half of the pairs.}
The factors from the second set of auxiliary primes have already been identified. The determinant coefficient also depends on what survives when the coordinates from the first set are specialized.

It remains to check the first-half determinant units at the
trivial character. For each signature \(s\), write \(H_{P,s}\) for its unramified Hecke polynomial with the prescribed character \(\eta_s\). Evaluation at one gives
\[
 H_{P,s}(1)=\det(1-\eta_s\rho_s(\mathrm{Frob}^{-1})).
\]

At the distinguished place the twist \(\eta_s\rho_s\)
acts by the scalar \(\delta\) or \(\delta^{-1}\) on inertia.
Intersect the kernels of all the reduced twisted matrices and
the conjugate of this intersection. This is an open normal
subgroup because there are finitely many continuous residual
actions. At either place of a first-half pair the inverse residual Frobenius is scalar. Its two possible values are
\[
 \overline{\delta(z)},\qquad \overline{\delta(z)}^{-1},
\]
since conjugation transports inertia and inverts \(\delta\). Thus the displayed determinant is a
unit for every signature.

Suppose \(\bar H_P(1)\) is singular over the algebraically
closed residue field. Its kernel is nonzero and stable under
all commuting Hecke operators. Restrict successively to a
nonzero eigenspace of a nonscalar operator. This gives a common
eigenvector \(w\) in the kernel. Its residual signature \(s\)
therefore satisfies
\[
       0=\bar H_P(1)w=\overline{H_{P,s}(1)}\,w .
\]
The first-half test just proved makes \(H_{P,s}(1)\) a unit,
a contradiction. Thus \(\det H_P(1)\) is a unit on the row.
Deleting the first-half pairs at the trivial character gives
\[
                   g_M(0)=\det C_0\prod_P\det H_P(1).
\]
This proves the coefficient and index bounds for every allowed
lattice. The following comparison identifies the integral lattice,
residual signature and the common geometric denominator.
\label{cor:twist-euler-2}
\end{proof}
